\chapter{Coronary Heart Disease}
\index{Coronary Heart Disease}

\section*{Definition and Overview}
Coronary heart disease (CHD), also known as ischaemic heart disease, is a condition in which the coronary arteries -- the blood vessels that supply the heart muscle -- become narrowed by the gradual build-up of fatty deposits. When these arteries cannot deliver enough blood and oxygen, the heart muscle is starved of its supply. This produces angina (chest pain or discomfort that typically occurs on exertion) and, if an artery becomes suddenly blocked by a blood clot, a heart attack (myocardial infarction). CHD is the most common form of heart disease in adults. It develops silently over many years, and many people have no symptoms until the narrowing is advanced.

\section*{Epidemiology}
Coronary heart disease is a leading cause of death and disability worldwide. In high-income countries it is, together with stroke, one of the two most common causes of death. The disease affects men at an earlier age than women, but after the menopause the risk in women rises to approach that of men. Risk increases steeply with age. CHD is more common in people with high blood pressure, high cholesterol, diabetes, or obesity, and in those who smoke or have a family history of early heart disease. While death rates have fallen in many developed countries thanks to better prevention and treatment, the burden remains high and is rising in low- and middle-income countries.

\section*{Aetiology and Pathophysiology}
\begin{itemize}
    \item \textbf{Atherosclerosis:} the underlying process in which fatty material, cholesterol, and inflammatory cells accumulate within the artery wall to form plaques
    \item \textbf{Plaque growth:} over years, plaques enlarge and narrow the coronary arteries, restricting blood flow to the heart muscle
    \item \textbf{Risk factors:} smoking, high blood pressure, high cholesterol, diabetes, obesity, physical inactivity, and a family history of premature heart disease
    \item \textbf{Angina mechanism:} when the demand for oxygen rises, as during exercise or emotional stress, narrowed arteries cannot deliver enough blood, causing chest discomfort that settles with rest
    \item \textbf{Heart attack mechanism:} a plaque can rupture, and a blood clot then forms on its surface, suddenly occluding the artery and damaging heart muscle
    \item \textbf{Ischaemia and infarction:} ischaemia is reduced blood flow to the heart muscle; if prolonged, it leads to irreversible muscle damage known as infarction
\end{itemize}

\section*{Clinical Features}
\begin{itemize}
    \item \textbf{Angina:} chest pain, pressure, heaviness, or tightness brought on by exertion or stress and relieved by rest or nitrates
    \item Pain may spread to the jaw, neck, back, shoulders, or either arm
    \item Breathlessness on exertion
    \item Fatigue and reduced exercise tolerance
    \item \textbf{Heart attack:} severe, crushing central chest pain lasting longer than a few minutes and not relieved by rest
    \item Heart attack pain may be accompanied by sweating, nausea, vomiting, and faintness
    \item Some people -- especially women, older adults, and people with diabetes -- present with atypical symptoms such as indigestion, back or jaw pain, or breathlessness alone
    \item \textbf{Silent ischaemia:} in some people, reduced blood flow produces no symptoms at all
\end{itemize}

\section*{Differential Diagnosis}
\begin{itemize}
    \item Gastro-oesophageal reflux and other causes of heartburn
    \item Costochondritis and other chest wall musculoskeletal pain
    \item Pericarditis, the inflammation of the sac surrounding the heart
    \item Aortic dissection -- a medical emergency causing sudden, tearing chest pain
    \item Pulmonary embolism, a blood clot in the lung
    \item Pneumonia and pleurisy
    \item Anxiety and panic attacks
    \item Abnormal heart rhythms (arrhythmias)
    \item Biliary disease, such as gallstones or cholecystitis
\end{itemize}

\section*{When to Seek Medical Care}
See a doctor if you develop new chest pain, chest tightness on exertion, unusual breathlessness, or any symptom that could be angina. Anyone with known CHD should be reviewed promptly if their angina becomes more frequent, easier to trigger, or begins to occur at rest.

\medskip
\textbf{Contact your local emergency services for:}
\begin{itemize}
    \item Chest pain or pressure lasting more than a few minutes, especially with sweating, nausea, or shortness of breath
    \item Pain spreading to the arm, jaw, neck, or back
    \item Sudden breathlessness, fainting, or collapse
    \item A rapid or irregular heartbeat that makes you feel faint
\end{itemize}

\section*{Diagnosis and Investigations}
\begin{itemize}
    \item \textbf{History:} careful description of the pain -- its nature, site, what triggers it, and what relieves it
    \item \textbf{Risk factor assessment:} measurement of blood pressure, cholesterol, blood glucose, plus smoking and family history
    \item \textbf{Electrocardiogram (ECG):} records the heart's electrical activity and may show changes of ischaemia or evidence of a previous heart attack
    \item \textbf{Blood tests:} cardiac troponin to detect heart muscle damage during a heart attack, together with cholesterol and glucose levels
    \item \textbf{Exercise stress testing:} an ECG recorded while walking on a treadmill to see whether the heart is starved of blood during exertion
    \item \textbf{Echocardiogram:} ultrasound assessment of heart pumping function, chamber size, and valve function
    \item \textbf{Coronary angiography:} an invasive procedure using dye and X-rays to visualise narrowed or blocked coronary arteries; it is the definitive test for blockages
    \item \textbf{CT coronary angiography:} a specialised CT scan that assesses calcium and plaque within the coronary arteries
    \item \textbf{Stress imaging:} nuclear or MRI scanning to identify regions of heart muscle not receiving enough blood during stress
\end{itemize}

\section*{Management}
\begin{itemize}
    \item \textbf{Lifestyle change:} stopping smoking, a heart-healthy diet, regular exercise, weight control, and limiting alcohol
    \item \textbf{Antiplatelet medicines:} low-dose aspirin or other agents to reduce the risk of blood clots
    \item \textbf{Statins:} cholesterol-lowering medicines that reduce LDL cholesterol and help stabilise plaques
    \item \textbf{Blood pressure control:} medicines to keep blood pressure within a healthy range
    \item \textbf{Antianginal therapy:} glyceryl trinitrate to relieve symptoms, and beta-blockers or calcium channel blockers to prevent them
    \item \textbf{Diabetes control:} careful blood-glucose management in people with diabetes
    \item \textbf{Cardiac rehabilitation:} a supervised program of exercise, education, and psychological support after a heart attack or procedure
    \item \textbf{Percutaneous coronary intervention:} angioplasty with stenting to open a narrowed or blocked artery
    \item \textbf{Coronary artery bypass grafting:} surgery using veins or arteries to bypass severe blockages
    \item \textbf{Follow-up:} regular monitoring of risk factors, symptoms, and medication adherence
\end{itemize}

\section*{Complications}
\begin{itemize}
    \item Heart attack (myocardial infarction) from sudden blockage of a coronary artery
    \item Heart failure, when damaged heart muscle can no longer pump effectively
    \item Abnormal heart rhythms, including life-threatening ventricular arrhythmias
    \item Cardiac arrest and sudden cardiac death
    \item Cardiogenic shock in very severe heart attacks
    \item Chronic angina and reduced quality of life
    \item Complications of treatment, including bleeding from antiplatelet drugs and stent or bypass graft problems
    \item Depression and anxiety after a heart attack
\end{itemize}

\section*{Prognosis and Outcomes}
With modern treatment, many people with coronary heart disease live for many years with a good quality of life. The outlook has improved substantially because of better medicines, rapid treatment of heart attacks, and the wider use of stents and bypass surgery. Outcomes are best when the disease is detected early, risk factors are well controlled, and the person adheres to their treatment plan. A heart attack remains a serious event, and recovery depends on how much muscle is damaged and how quickly treatment is delivered. Long-term outlook is strongly influenced by stopping smoking, taking prescribed medicines reliably, and maintaining a healthy lifestyle. Individual prognosis should be discussed with the treating cardiologist.

\section*{Prevention and Self-Care}
\begin{itemize}
    \item Do not smoke, and seek help to quit if needed
    \item Eat a heart-healthy diet rich in vegetables, fruit, whole grains, and fish, and low in saturated fat and salt
    \item Stay physically active, aiming for at least 150 minutes of moderate activity per week
    \item Maintain a healthy weight and waist measurement
    \item Limit alcohol to within recommended guidelines
    \item Keep blood pressure, cholesterol, and blood glucose within target ranges, taking prescribed medicines
    \item Attend regular health checks and cardiovascular risk screening
    \item Manage stress and seek help for anxiety or low mood
    \item Learn the symptoms of a heart attack and make an emergency plan
    \item Follow cardiac rehabilitation advice after a heart attack or procedure
\end{itemize}

\section*{Sources and Further Reading}
The following reputable sources provide further information on coronary heart disease:
\begin{itemize}
    \item NHS. \textit{Health A--Z: coronary heart disease information.} https://www.nhs.uk/conditions/
    \item Mayo Clinic. \textit{Diseases and conditions library.} https://www.mayoclinic.org/
    \item British Heart Foundation. \textit{Heart conditions and treatments information.} https://www.bhf.org.uk/
    \item NICE. \textit{Clinical guidelines on cardiovascular disease.} https://www.nice.org.uk/
    \item MSD Manual. \textit{Professional edition: cardiovascular disorders.} https://www.msdmanuals.com/
    \item World Health Organization. \textit{Cardiovascular disease fact sheets.} https://www.who.int/
\end{itemize}
